\clearpage
\begin{proofbox}
  \textbf{\textcolor{CProof}{方法一：把 $y$ 当成已知数，完整证明通用结论（最重要）}}
\begin{center}\begin{tikzpicture}[font=\small,>=stealth,line join=round]
\draw[gray!40,rounded corners=3pt] (0,0) rectangle (14.2,9.6);
\node at (7.1,9.0) {判别式筛出候选范围；二次项消失时，另查方程};
\pgfmathsetmacro{\LL}{(3-sqrt(5))/2}
\pgfmathsetmacro{\UU}{(3+sqrt(5))/2}
\pgfmathsetmacro{\rr}{(sqrt(5)-1)/2}
\pgfmathsetmacro{\ss}{-(sqrt(5)+1)/2}
\begin{scope}[shift={(2.0,3.0)},x=1.6cm,y=.70cm]
\draw[->,gray] (-.2,0)--(3.5,0) node[right]{$y$};
\draw[->,gray] (0,-1.5)--(0,6) node[above]{$\Delta$};
\fill[blue!8,domain=\LL:\UU,samples=100] (\LL,0)--plot(\x,{4*(-\x*\x+3*\x-1)})--(\UU,0)--cycle;
\draw[teal!80!black,very thick,domain=.2:2.8,samples=120] plot(\x,{4*(-\x*\x+3*\x-1)});
\draw[blue,very thick] (\LL,0)--(\UU,0);
\draw[dotted,gray] (\LL,0)--(\LL,-2);\draw[dotted,gray] (\UU,0)--(\UU,-2);\node[below=5pt] at (\LL,-2){$L$};\node[below=5pt] at (\UU,-2){$U$};
\draw[dashed,orange!85!black] (1,0)--(1,4);
\fill[orange!85!black] (1,4) circle(2pt);\node[below=5pt] at (1,0){$1$};
\end{scope}
\node[align=left,anchor=west] at (8.2,5.1){$\Delta(y)=4(-y^2+3y-1)$\\[12pt]$\Delta\ge0\iff L\le y\le U$\\[12pt]橙点对应 $y=1$\\[6pt]此时改查 $-2x-1=0$\\[6pt]$x=-1/2$，仍能取到};
\node[text width=13cm,align=center] at (7.1,.6) {横轴是输出值 $y$；这不是原函数关于 $x$ 的图像。};
\end{tikzpicture}
\end{center}


  \medskip
  \textbf{第 1 步：先确认分母能不能为 $0$。}
  设 $Q(x)=dx^2+ex+f$，其中 $d\ne0$，并已知 $e^2-4df<0$。
  由
  \[
    Q(x)=d\left(x+\frac{e}{2d}\right)^2+\frac{4df-e^2}{4d}
  \]
  可知 $Q(x)$ 始终与 $d$ 同号，所以对任意实数 $x$，$Q(x)\ne0$。
  这样两边同乘 $Q(x)$ 不会凭空添出不允许的 $x$。

  \medskip
  \textbf{第 2 步：把函数变成关于 $x$ 的方程。}
  \[
    y=\frac{ax^2+bx+c}{dx^2+ex+f}
    \quad\Longleftrightarrow\quad
    y(dx^2+ex+f)=ax^2+bx+c,
  \]
  移项，得到
  \[
    \boxed{(a-dy)x^2+(b-ey)x+(c-fy)=0.}\tag{1}
  \]
  这一步的意思是：指定一个 $y$ 后，就去找能让它成立的 $x$。

  \medskip
  \textbf{第 3 步：先处理真正的二次方程。}
  当 $a-dy\ne0$ 时，令
  \[
    A=a-dy,\qquad B=b-ey,\qquad C=c-fy.
  \]
  方程 (1) 就是 $Ax^2+Bx+C=0$。求根公式给出
  \[
    x=\frac{-B\pm\sqrt{B^2-4AC}}{2A}.
  \]
  因为 $A\ne0$，要找到实数 $x$，就必须且只需
  \[
    B^2-4AC\ge0,
  \]
  也就是
  \[
    \boxed{(b-ey)^2-4(a-dy)(c-fy)\ge0.}\tag{2}
  \]
  这既说明\textbf{为什么必须满足这个不等式}，也说明\textbf{满足它就一定能找到实数 $x$}。

  \medskip
  \textbf{第 4 步：整理成只含 $y$ 的不等式。}
  \[
    \begin{aligned}
      \Delta(y)
      &=(b-ey)^2-4(a-dy)(c-fy)\\
      &=b^2-2bey+e^2y^2\\
      &\quad-4\bigl(ac-afy-cdy+dfy^2\bigr)\\
      &=(e^2-4df)y^2\\
      &\quad+(4af+4cd-2be)y+(b^2-4ac).
    \end{aligned}
  \]
  解 $\Delta(y)\ge0$，得到候选范围；这里二次项 $e^2-4df<0$，因此对应的是二次式开口向下时的中间一段。

  \medskip
  \textbf{第 5 步：单独检查 $y=\dfrac ad$。}
  当 $y=\dfrac ad$ 时，方程 (1) 不再含 $x^2$，而变成
  \[
    \left(b-\frac{ae}{d}\right)x+\left(c-\frac{af}{d}\right)=0.
  \]
  把两边乘以 $d$，得到更容易判断的一次方程：
  \[
    \boxed{(bd-ae)x+(cd-af)=0.}\tag{3}
  \]
\end{proofbox}
\clearpage
\begin{proofbox}
用三个具体函数，看剩余方程的三种不同情况：
\begin{center}\begin{tikzpicture}[font=\small,>=stealth,line join=round]
\draw[gray!40,rounded corners=3pt] (0,0) rectangle (14.2,13.5);
\node at (7.1,12.9) {二次项消失之后：改看剩余式子 $H(x)$ 有没有零点};
\begin{scope}[shift={(3.6,10.3)},x=1cm,y=.8cm]
\draw[->,gray] (-2.2,0)--(2.3,0) node[right]{$x$};
\draw[->,gray] (0,-1.5)--(0,1.65) node[above]{$H$};
\draw[blue,very thick,domain=-1.1:.1] plot(\x,{-2*\x-1});\fill[teal!80!black] (-.5,0) circle(2.2pt);
\end{scope}
\node[align=left,anchor=west] at (7.3,10.45) {$y=1$：$-2x-1=0$\\[8pt]与横轴相交，$x=-1/2$ 可用};
\begin{scope}[shift={(3.6,6.6)},x=1cm,y=.8cm]
\draw[->,gray] (-2.2,0)--(2.3,0) node[right]{$x$};
\draw[->,gray] (0,-1.5)--(0,1.65) node[above]{$H$};
\draw[blue,very thick] (-2,1)--(2,1);
\end{scope}
\node[align=left,anchor=west] at (7.3,6.75) {$y=2$：$1=0$\\[8pt]始终离开横轴，无解};
\begin{scope}[shift={(3.6,2.9)},x=1cm,y=.8cm]
\draw[->,gray] (-2.2,0)--(2.3,0) node[right]{$x$};
\draw[->,gray] (0,-1.5)--(0,1.65) node[above]{$H$};
\draw[blue,very thick] (-2,0)--(2,0);
\end{scope}
\node[align=left,anchor=west] at (7.3,3.05) {$y=2$：$0=0$\\[8pt]整条横轴都是零点，任意 $x$ 可用};

\node[text width=13cm,align=center] at (7.1,.6) {从上到下对应标准题、缺端点题、常函数题；仍须确认原分母不为零。};
\end{tikzpicture}
\end{center}
一般地，按系数分成三种情况：
  \begin{itemize}[leftmargin=*,itemsep=0.35em]
    \item 若 $bd-ae\ne0$，可取 $x=\dfrac{af-cd}{bd-ae}$，所以 $y=\dfrac ad$ \textbf{能取到}。
    \item 若 $bd-ae=0$ 且 $cd-af\ne0$，方程变成一个非零常数等于 $0$，不可能成立，所以 $y=\dfrac ad$ \textbf{取不到}。
    \item 若 $bd-ae=0$ 且 $cd-af=0$，方程变成 $0=0$，任意 $x$ 都可以；这时整个函数恒等于 $\dfrac ad$，值域只有这一个数。
  \end{itemize}
  \textbf{特别注意：}即使把 $y=\dfrac ad$ 代入 (2) 得到 $\Delta(y)=0$，也\textbf{不能因此就认定它能取到}；因为此时已不是二次方程。

  \medskip
  \textbf{第 6 步：合并答案。}
  保留满足 (2) 且 $a-dy\ne0$ 的所有 $y$，再按照第 5 步决定是否加入 $\dfrac ad$。
  这样既没有漏掉真正能取到的值，也不会把不能取到的值算进去。

  \bigskip
  \end{proofbox}
\clearpage
\begin{proofbox}
\textbf{\textcolor{CProof}{方法二：先拆开分式，用 $x^2\ge0$ 求范围（初中也能理解）}}
\begin{center}\begin{tikzpicture}[font=\small,>=stealth,line join=round]
\draw[gray!45,rounded corners=3pt] (0,0) rectangle (14.2,9.4);
\node at (7.1,8.8) {函数无限接近下边界，但从不落到下边界上};

\begin{scope}[shift={(4.1,1.8)},x=.9cm,y=1.65cm]
\draw[->] (-3.5,0)--(3.6,0) node[right]{$x$};\draw[->] (0,-.1)--(0,3.8) node[above]{$y$};\draw[dashed,magenta] (-3.3,2)--(3.3,2);
\draw[teal!80!black,thick,domain=-3.2:3.2,samples=100] plot(\x,{2+1/(\x*\x+1)});\fill[blue] (0,3) circle(2pt);\node[above right] at (0,3){$(0,3)$};\node[above left] at (0,2){$2$};
\end{scope}
\node[align=left,anchor=west] at (8.3,5.1){$y=2+\dfrac1{x^2+1}$\\[10pt] $0<\dfrac1{x^2+1}\le1$\\[10pt] 值域 $(2,3]$\\[12pt] $y=2$ 时变成 $1=0$};
\node[text width=13.2cm,align=center] at (7.1,.65) {即使候选判别式在 $y=2$ 处为零，降次后也可能无解。};
\end{tikzpicture}
\end{center}


  \medskip
  下面用一个典型题说明：
  \[
    y=\frac{2x^2+3}{x^2+1}.
  \]
  直接把分子写成 $2(x^2+1)+1$：
  \[
    y=2+\frac1{x^2+1}.
  \]
  因为 $x^2\ge0$，所以 $x^2+1\ge1$；同时分母始终是正数，故
  \[
    0<\frac1{x^2+1}\le1.
  \]
  两边同时加 $2$，就得到
  \[
    \boxed{2<y\le3.}
  \]
  $x=0$ 时 $y=3$，所以右端点能取到；无论 $x$ 取什么实数，$\dfrac1{x^2+1}$ 都大于 $0$，因此左端点 $2$ 永远取不到。
  这道题\textbf{拆分法比判别式更快}。并且它也解释了为什么“越来越接近 $2$”不等于“能等于 $2$”。

  \bigskip
  \end{proofbox}
\clearpage
\begin{proofbox}
\textbf{\textcolor{CProof}{方法三：利用平方不小于 $0$，同时锁定最大值和最小值}}
\begin{center}\begin{tikzpicture}[font=\small,>=stealth,line join=round]
\draw[gray!45,rounded corners=3pt] (0,0) rectangle (14.2,9.4);
\node at (7.1,8.8) {配平方：与上下边界的差，都不可能小于零};

\begin{scope}[shift={(4.3,2.3)},x=.85cm,y=1.5cm]
\draw[->] (-4,0)--(3.3,0) node[right]{$x$};\draw[->] (0,-.2)--(0,3) node[above]{差值};
\draw[blue,thick,domain=-3.8:3.1,samples=120] plot(\x,{(\x*\x+1)/(\x*\x+2*\x+2)-.38196601125});
\draw[orange!90!black,thick,domain=-3.8:3.1,samples=120] plot(\x,{2.61803398875-(\x*\x+1)/(\x*\x+2*\x+2)});
\fill[blue] (.618034,0) circle(2pt);\fill[orange!90!black] (-1.618034,0) circle(2pt);\node[below] at (.618034,0){$r_1$};\node[below] at (-1.618034,0){$r_2$};
\end{scope}
\node[align=left,anchor=west] at (8.4,5.2){\textcolor{blue}{$F-L\ge0$}，$x=r_1$ 取零\\[10pt] \textcolor{orange!90!black}{$U-F\ge0$}，$x=r_2$ 取零\\[12pt] $r_1=(\sqrt5-1)/2$\\[8pt] $r_2=-(\sqrt5+1)/2$};
\node[text width=13.2cm,align=center] at (7.1,.65) {图画的是两个差值；分子配平方、分母恒正，保证两条曲线均不穿过横轴。};
\end{tikzpicture}
\end{center}


  \medskip
  对稍复杂的题
  \[
    F(x)=\frac{x^2+1}{x^2+2x+2},\qquad
    x^2+2x+2=(x+1)^2+1>0,
  \]
  我们想证明它的值始终夹在两个数之间。
  设
  \[
    L=\frac{3-\sqrt5}{2},\qquad
    U=\frac{3+\sqrt5}{2}.
  \]
  先把 $F(x)-L$ 通分，再把分子配成平方：
  \[
    \begin{aligned}
      F(x)-L
      &=\frac{(1-L)x^2-2Lx+(1-2L)}{x^2+2x+2}\\
      &=\frac{(1-L)\left(x-\dfrac{\sqrt5-1}{2}\right)^2}
      {x^2+2x+2}\ge0.
    \end{aligned}
  \]
  因为 $1-L>0$，所以上式确实不小于 $0$，即 $F(x)\ge L$。
  当 $x=\dfrac{\sqrt5-1}{2}$ 时，平方恰好是 $0$，因此最小值 $L$ 真能取到。

  \medskip
  同样地，
  \[
    \begin{aligned}
      U-F(x)
      &=\frac{(U-1)x^2+2Ux+(2U-1)}{x^2+2x+2}\\
      &=\frac{(U-1)\left(x+\dfrac{\sqrt5+1}{2}\right)^2}
      {x^2+2x+2}\ge0.
    \end{aligned}
  \]
  所以 $F(x)\le U$。当 $x=-\dfrac{\sqrt5+1}{2}$ 时等号成立，最大值 $U$ 也确实能取到。
  因此
  \[
    \boxed{\frac{3-\sqrt5}{2}\le F(x)\le\frac{3+\sqrt5}{2}.}
  \]
  \textbf{为什么中间不会缺值？}因为分母始终不为 $0$，函数图像不会中途断开；从取得最小值的那个 $x$ 连续移动到取得最大值的那个 $x$，中间的高度也都会经过。

  \medskip
  \textbf{这两个数从哪里来？}它们并不是靠猜的：用方法一列出
  \[
    (1-y)x^2-2yx+(1-2y)=0
  \]
  并令 $\Delta(y)=0$，就得到 $y^2-3y+1=0$，解出的两个数恰好是 $L$ 与 $U$。
  方法一善于\textbf{找到边界}，方法三则能用平方\textbf{直接解释为什么越不过边界、什么时候取到边界}。

  \medskip
  \textbf{\textcolor{CProof}{结论回扣}}
  三种方法不是每题都要用：\textbf{能拆就拆，拆不开用判别式；想进一步解释最值为什么成立，可以配平方验证。}
\end{proofbox}


% ==== 新增：直接用原来的 x 的两种数形结合证明，不删除原三种证明 ====
\clearpage
\begin{proofbox}
  \textbf{\textcolor{CProof}{几何补充一：抛物线斜率法（从分式本身自然想到）}}

  \medskip
  \textbf{问题：}不用先求导、不改变原来的自变量，求
  \[
    F(x)=\frac{x^2+1}{x^2+2x+2}
  \]
  的值域，并用图形理解判别式等于 $0$。

  \textbf{第 1 步：把相除改看成斜率。}
  设
  \[
    X=x^2+2x+2,\qquad Y=x^2+1.
  \]
  $X=(x+1)^2+1>0$，因此 $F(x)=Y/X$，就是原点与 $M(X,Y)$ 连线的斜率。
  这里的 $X,Y$ 只是点的两个坐标，\textbf{原自变量始终还是 $x$}。

  \textbf{第 2 步：找到点 $M$ 的运动轨迹。}
  \[
    X-Y=2x+1\quad\Longrightarrow\quad x=\frac{X-Y-1}{2}.
  \]
  又因为 $Y-1=x^2$，得到
  \[
    \boxed{(X-Y-1)^2=4(Y-1).}
  \]
  这是一条转过方向的抛物线。下面的图中，两条从原点出发的边界直线与它相切。
  \begin{center}\begin{tikzpicture}[font=\small,>=stealth,line join=round]
\draw[gray!40,rounded corners=3pt] (0,0) rectangle (14.2,9.8);
\node at (7.1,9.200000000000001) {两条原点切线夹住轨迹：边界函数值就是边界斜率};
\pgfmathsetmacro{\LL}{(3-sqrt(5))/2}
\pgfmathsetmacro{\UU}{(3+sqrt(5))/2}
\pgfmathsetmacro{\rr}{(sqrt(5)-1)/2}
\pgfmathsetmacro{\ss}{-(sqrt(5)+1)/2}
\begin{scope}[shift={(1.3,1.8)},x=.48cm,y=.48cm]
\draw[->,gray] (0,0)--(13,0) node[right]{$X$};
\draw[->,gray] (0,0)--(0,12) node[above]{$Y$};
\draw[blue,very thick,domain=-3.1:2.2,samples=180,variable=\t] plot ({\t*\t+2*\t+2},{\t*\t+1});
\draw[orange!85!black,thick] (0,0)--(12,{\LL*12});
\draw[teal!80!black,thick] (0,0)--(4.2,{\UU*4.2});
\fill[orange!85!black] ({(5+sqrt(5))/2},{(5-sqrt(5))/2}) circle(2pt) node[below right=5pt]{$A$};
\fill[teal!80!black] ({(5-sqrt(5))/2},{(5+sqrt(5))/2}) circle(2pt) node[above left=5pt]{$B$};
\node[below left=4pt] at (0,0){$O$};
\end{scope}
\node[align=left,anchor=west] at (8.1,5.4){$X=x^2+2x+2>0$\\[8pt]$Y=x^2+1$\\[12pt]$F(x)=Y/X$ 是连线斜率\\[12pt]\textcolor{orange!85!black}{$OA$ 斜率 $L=(3-\sqrt5)/2$}\\[10pt]\textcolor{teal!80!black}{$OB$ 斜率 $U=(3+\sqrt5)/2$}};
\node[text width=13cm,align=center] at (7.1,.6) {坐标轴等比例；$A,B$ 均来自有限实数 $x$，两条边界斜率都能取到。};
\end{tikzpicture}
\end{center}

  \textbf{第 3 步：用直线与曲线有无交点求斜率。}
  设从原点出发的直线为 $Y=kX$。
  代入最初的 $X,Y$，得到
  \[
    x^2+1=k(x^2+2x+2)
    \quad\Longleftrightarrow\quad
    (1-k)x^2-2kx+(1-2k)=0. \tag{4}
  \]
  对于 $k\ne1$，它是关于 $x$ 的二次方程；有实数交点的条件为
  \[
    \begin{aligned}
      \Delta(k)&=(-2k)^2-4(1-k)(1-2k)\\
      &=4(-k^2+3k-1)\ge0.
    \end{aligned}
  \]
  所以
  \[
    \boxed{\frac{3-\sqrt5}{2}\le k\le\frac{3+\sqrt5}{2}.}
  \]
  $k=1$ 时，方程 (4) 变成 $-2x-1=0$，解得 $x=-\frac12$；这个斜率也能取到。

  \textbf{第 4 步：说明两个端点都确实能取到。}
  两条边界直线分别对应 $\Delta(k)=0$，即直线与抛物线相切。
  当 $k=(3-\sqrt5)/2$ 时，$x=(\sqrt5-1)/2$；
  当 $k=(3+\sqrt5)/2$ 时，$x=-(\sqrt5+1)/2$。
  两个数都是真实可取的 $x$，分母也都不为 $0$，所以
  \[
    \boxed{F(x)\in\left[\frac{3-\sqrt5}{2},\frac{3+\sqrt5}{2}\right].}
  \]

  \textbf{几何意义：}对这道题而言，$\Delta>0$ 对应两处相交，$\Delta=0$ 对应相切，$\Delta<0$ 对应不相交；
  但在 $k=1$ 这个二次项消失的位置，只有一个有限交点却\textbf{不是相切}，必须单独判断。
\end{proofbox}

\clearpage
\begin{proofbox}
  \textbf{\textcolor{CProof}{几何补充二：单位圆法（从平方差自然找到勾股关系）}}

  \medskip
  还是研究同一个函数，完全保留原来的 $x$：
  \[
    F(x)=\frac{x^2+1}{D},\qquad D=x^2+2x+2>0.
  \]
  \textbf{第 1 步：问分母能否变成直角三角形的斜边。}
  因为 $D-1=x^2+2x+1$ 是平方，而
  \[
    D^2-(D-2)^2=4(D-1)=(2x+2)^2,
  \]
  所以
  \[
    \boxed{(D-2)^2+(2x+2)^2=D^2.}\tag{5}
  \]
  这就是勾股关系的形式。$D-2$、$2x+2$ 是有符号的量，作为边长时应取绝对值；下面选两者都为正的 $x=1$，先看一个真正的直角三角形。
\begin{center}\begin{tikzpicture}[font=\small,>=stealth,line join=round]
\draw[gray!40,rounded corners=3pt] (0,0) rectangle (14.2,8.6);
\node at (7.1,8.0) {先看一个真正的直角三角形，再把平方关系归一化};
\coordinate (A) at (1.7,2);\coordinate (B) at (5.3,2);\coordinate (C) at (5.3,6.8);
\draw[blue,very thick] (A)--(B)--(C)--cycle;
\draw[gray] (4.95,2)--(4.95,2.35)--(5.3,2.35);
\node[below=7pt] at (3.5,2){$D-2=3$};
\node[right=7pt] at (5.3,4.4){$2x+2=4$};
\node[above left=5pt] at (3.5,4.4){$D=5$};
\node[align=left,anchor=west] at (8.2,5.3){取 $x=1$，则 $D=5$\\[12pt]$3^2+4^2=5^2$\\[12pt]同除以 $D^2$：\\[8pt]$(3/5)^2+(4/5)^2=1$\\[12pt]圆上点为 $(3/5,4/5)$};
\node[text width=13cm,align=center] at (7.1,.6) {本图选 $x=1$，两直角边为正；一般情形的边长应取 $|D-2|$、$|2x+2|$。};
\end{tikzpicture}
\end{center}


  \textbf{第 2 步：把斜边缩成 $1$，单位圆便出现了。}
  令
  \[
    u=\frac{D-2}{D},\qquad v=\frac{2x+2}{D}.
  \]
  (5) 两边除以 $D^2$，得到
  \[
    \boxed{u^2+v^2=1.}
  \]
  这里的 $u,v$ 是\textbf{圆上点的坐标}，不是在更换原来的自变量。

  \textbf{第 3 步：让原函数变成圆上点的一个简单式子。}
  由 $u=1-2/D$ 可知 $1/D=(1-u)/2$；又有
  \[
    x^2+1=D-(2x+1)=D-(2x+2)+1.
  \]
  两边除以 $D$，可得
  \[
    F(x)=1-v+\frac1D
    =\boxed{\frac32-\frac12u-v.}
  \]
  于是每给定一个函数值 $y$，就对应一条直线
  \[
    \frac12u+v=\frac32-y.
  \]
  要让该直线和半径为 $1$ 的圆相交，它到圆心的距离不能超过 $1$：
  \[
    \frac{|\frac32-y|}{\sqrt{(\frac12)^2+1}}\le1
    \quad\Longleftrightarrow\quad
    \left|\frac32-y\right|\le\frac{\sqrt5}{2}.
  \]
  因而先得到候选范围
  \[
    \frac{3-\sqrt5}{2}\le y\le\frac{3+\sqrt5}{2}.
  \]
  \begin{center}\begin{tikzpicture}[font=\small,>=stealth,line join=round]
\draw[gray!40,rounded corners=3pt] (0,0) rectangle (14.2,10.2);
\node at (7.1,9.6) {单位圆与平行直线相切：两个边界都避开缺点 $E$};
\begin{scope}[shift={(4.0,4.6)},x=2cm,y=2cm]
\draw[->,gray] (-1.3,0)--(1.45,0) node[right]{$u$};
\draw[->,gray] (0,-1.4)--(0,1.6) node[above]{$v$};
\draw[blue,very thick] (0,0) circle(1);
\draw[teal!80!black,thick] (-1.3,{sqrt(5)/2+.65})--(1.3,{sqrt(5)/2-.65});
\draw[orange!85!black,thick] (-1.3,{-sqrt(5)/2+.65})--(1.3,{-sqrt(5)/2-.65});
\fill[teal!80!black] ({1/sqrt(5)},{2/sqrt(5)}) circle(2pt) node[above right=5pt]{$A$};
\fill[orange!85!black] ({-1/sqrt(5)},{-2/sqrt(5)}) circle(2pt) node[below left=5pt]{$B$};
\draw[red!70!black,thick,fill=white] (1,0) circle(2.8pt) node[below right=6pt]{$E$};
\node[below left=4pt] at (0,0){$O$};
\end{scope}
\node[align=left,anchor=west] at (8.1,5.6){$u^2+v^2=1$\\[10pt]$\tfrac12u+v=\tfrac32-y$\\[12pt]\textcolor{teal!80!black}{$A$ 对应最小值 $L$}\\[9pt]\textcolor{orange!85!black}{$B$ 对应最大值 $U$}\\[12pt]\textcolor{red!70!black}{$E=(1,0)$ 不可达}\\[7pt]两切点都不是 $E$};
\node[text width=13cm,align=center] at (7.1,.6) {圆的横纵坐标等比例；空心点不属于原参数能够到达的轨迹。};
\end{tikzpicture}
\end{center}

  \textbf{第 4 步：圆并不是每一个点都能取到，端点必须查。}
  对任何有限的实数 $x$，
  \[
    1-u=\frac2D>0,
  \]
  所以 $(1,0)$ 是这个圆上\textbf{唯一不能到达的点}。
  反过来，只要 $u\ne1$，就能用
  \[
    x=\frac{v}{1-u}-1
  \]
  找回原来的实数 $x$。
  本题两条边界直线与圆相切的点分别为
  \[
    \left(\frac1{\sqrt5},\frac2{\sqrt5}\right),\qquad
    \left(-\frac1{\sqrt5},-\frac2{\sqrt5}\right).
  \]
  它们都不是 $(1,0)$，对应的 $x$ 也分别为
  \[
    \frac{\sqrt5-1}{2},\qquad-\frac{\sqrt5+1}{2}.
  \]
  因此本题上下边界\textbf{都能取到}，值域就是
  \[
    \boxed{\left[\frac{3-\sqrt5}{2},\frac{3+\sqrt5}{2}\right].}
  \]

  \textbf{两种图形的关系：}抛物线来自“分式等于斜率”，最自然地解释判别式；
  单位圆来自“分母里的平方差”，适合进一步研究相切与端点取舍。
  \textbf{不能为了用圆而先凭空设两个分式，也不能忽略圆上的缺口。}
\end{proofbox}

